\begin{textsample}{NEG Dim 5 – vem\_ed\_10 – Score 1.00 – vem\_ed\_10/t039.txt}  \label{ex:f5_neg_037}
QUEBRA-GELO Osvaldo Succi Junior Coordenador dos PCIs Em 26 de novembro de 2021, realizamos a Jornada de Projetos Colaborativos Internacionais (PCIs / Cesu): Reflexões sobre as Práticas, para discutir algumas das melhores práticas utilizadas nas Fatecs na realização de Intercâmbios Virtuais (ou, como são conhecidos no Centro Paula Souza, PCIs).
Na realização do evento online, tivemos o apoio da Fatec Guaratinguetá e a participação das Fatecs Americana, Baixada Santista, Bauru, Garça, Guaratinguetá, Itaquaquecetuba, Itu, Lins, Piracicaba, São José do Rio Preto e São Caetano do Sul.
Esta edição é dedicada inteiramente à Jornada, que teve quatro painéis temáticos: “Gestão da Internacionalização em Casa”, “Diversidade em Projetos”, “We are SOME of the champions” (relatos de quatro professores com experiência consolidada em PCIs) e “Pesquisas em PCIs”.
Rafael Ferreira Alves, coordenador da Unidade do Ensino Superior de Graduação / Cesu, afirma que a “abordagem de desenvolvimento dos PCIs incorpora metodologias ativas e situações desafiadoras, com interações colaborativas entre docentes e discentes de diferentes países.
Agrega ao perfil do egresso competências profissionais e socioemocionais únicas, mantendo aderência ao currículo do curso e fortalecendo o processo de ensino-aprendizagem na formação profissional e tecnológica.
A modelagem dos PCIs se relaciona literalmente à indissociabilidade do Ensino, Pesquisa e Extensão, cujo conceito é considerado nos indicadores de qualidade do Sistema Nacional de \textbf{Avaliação} da Educação Superior (Sinaes).”

% matched lemmas: avaliação
\end{textsample}
